
%%%%%%%%%%%%%%%%%%%%%%%%%%%%%%%%%%%%%%%%%%%%%%%%%%%%%%%%%%%%%
\section{模型的评价}

\subsection{模型的优点}
\begin{itemize}[itemindent=2em]
\item \textbf{数学结构清晰、可解释性强。} Prophet将需求分解为趋势、季节性、节假日等可命名和可视化的分量，报童模型是易腐品库存管理的标准框架，加成率、订货量、折扣回收率等参数均可直接对应商超的经营决策变量。
\item \textbf{方法选择有数据支撑。} Spearman替代Pearson基于全部六个品类正态性检验失败的数据事实（$p<10^{-75}$）；Prophet替代ARIMA基于留出集RMSE的系统性优势；自有弹性估计替代文献值的决策已有数据验证（与聂森等差异达8倍）。
\item \textbf{基线可比、增益可分解。} Q2和Q3的Baseline均使用与主方法完全相同的利润函数、参数和数据，增益可分解为预测精度提升和定价优化两部分，分析透明。
\item \textbf{稳健性透明。} 关键参数敏感性全部量化：$k$高敏感（$\pm 20\%$），弹性、损耗率、订货量上限均为低敏感，$k$的敏感性已在论文中以区间值形式报告。
\item \textbf{递进式建模。} Q1的描述性分析为Q2提供特征选择依据，Q2的品类级框架为Q3的单品级优化提供需求基准和弹性输入，四问逻辑连贯。
\end{itemize}

\subsection{模型的缺点}
\begin{itemize}[itemindent=2em]
\item \textbf{需求识别偏误。} 无法区分零销量=无需求与零销量=缺货，所有需求估计实质上是对实际销售量而非真实需求量的建模。87个稀疏单品的需求估计可靠性低。
\item \textbf{预测外推风险。} Prophet在2023年6月留出集上表现良好，但7月1--7日为纯预测，无验证数据。疫情后消费模式变化可能降低历史数据代表性。
\item \textbf{简单弹性估计。} Double-log回归$R^2$仅0.011--0.144，遗漏变量（天气、竞争、促销）可能影响弹性精度。但稳健性分析表明弹性误差对最优解影响极小。
\item \textbf{静态损耗率与单周期假设。} 损耗率假设恒定，未考虑季节性损耗差异。每日独立运行报童优化，未考虑日间需求序列相关。
\item \textbf{Q3选品无替代/互补关系。} 过滤法独立判断每个单品是否达标，不考虑单品间的需求替代或互补效应。
\end{itemize}
